\paragraph{Four different sizes.}
Suppose four nonzero conjugates have increasing absolute values $r_i$ with $r_{i+1}/r_i\ge16$. For exponents $e_1<\cdots<e_4$, multiply column $j$ of $(z_i^{e_j})$ by a positive $c_j$ with $c_{j+1}/c_j=(r_jr_{j+1})^{-(e_{j+1}-e_j)/2}$. On either side of the diagonal the absolute entries decrease by factors of at least four as one moves away from it. The sum of the other entries in each row is at most $2/3$ of the diagonal entry in absolute value. Thus the matrix is nonsingular. Any relation with fewer terms could be padded by unused exponents. For tighter gaps the earlier checker instead sums all 23 nondiagonal permutation magnitudes relative to the diagonal. The sum is largest at unit exponent gaps, by rearrangement, and a bound strictly below one suffices.

\paragraph{The analytic cutoff.}
Put $D=|\Delta|$ and $R=\sqrt D$. A split prime $p\nmid D$, $p\le R/2$, gives a primitive reduced form with leading coefficient $p$. Indeed one chooses $b^2\equiv\Delta\pmod{4p}$, $|b|\le p$, and puts $c=(b^2+D)/(4p)\ge p$. The absolute value of its singular modulus is within $2079$ of $\exp(\pi R/p)$ by the cusp estimate in \cite{riffaut}. The principal form gives leading coefficient 1. For $D\ge10^{11}$ three eligible primes occur, respectively, in $(0,.046R]$, $(.048R,.189R]$, $(.230R,.500R]$. To record the constants, set
\[
 \begin{split}
 E(y)&=\sqrt y\big[(\log y/(2\pi)+2)\log D+2(\log^2 y/(8\pi)+2)\big],\\
 B(y)&=2\sqrt y\log y+\log D+2\log 2,\qquad E(0)=0 .
 \end{split}
\]
Greni\'e and Molteni's prime ideal estimate for $K$ gives $|\psi_K(y)-y|\le E(y)$ for $y\ge100$. Its hypothesis follows from Dirichlet GRH. Over rational primes $p\le\sqrt y$, prime ideal powers contribute at most $2\sqrt y\log y$. The ramified primes and the split primes dividing $f$ cost at most $\log|d_K|+2\log f=\log D$ at larger primes. The extra term allows removal of $2$. Thus twice the eligible prime log-weight on $(x,y]$ is at least $y-x-E(x)-E(y)-B(y)$. At the three intervals the certified lower margins are respectively $68,191,751$ at $D=10^{11}$. After division by $R$ the error terms decrease: each monomial in $E(cR)/R$ contains $R^{-1/2}$ and at most two logarithm factors, each at least 9 for the positive endpoints. The same monotonicity holds for $B(cR)/R$. For two consecutive prime intervals with adjacent endpoints $vR,u'R$, the cusp ratio is bounded below by $(e^{\pi/v}-2079)/(e^{\pi/u'}+2079)>16$; the principal-to-first ratio is even larger since $p\ge3$. This proves the large range.\cite{gm,prior}

\paragraph{The finite range.}
For $10^7\le D<10^{11}$, enumerate entire residue classes where $1,p,p^2,p^3$ give four eligible reduced forms for a split prime $p$ with $4p^6\le D$. Include the powers of 2 when $D\equiv7\pmod 8$. Lifting solutions of $b^2\equiv\Delta\pmod{4p^j}$ produces the forms. For consecutive leading coefficients $a<b$, the tests $D\ge16a^2$ and $9D(b-a)^2\ge16a^2b^2$ imply a ratio exceeding 16 by the cusp bound. They hold for these powers in the specified range. The rest are tested individually. The CRT sieve accounts for all $49{,}995{,}000{,}000$ discriminants in the half-open range, including $47{,}116{,}196$ settled with three individually tested split primes. Each bin balances against the exact count in residue classes 0 and 3 modulo 4.

Below $10^7$, the integer form sieve and interval determinant test leave 80 inputs of degree at least four. The reduction enumerates primitive forms with $0\le b\le a\le c$ for absolute values, treating nonreal conjugate pairs once. The published output also covers the endpoint $10^7$. There are 67 low-degree inputs, $4{,}999{,}603$ exclusions at the first pass and another 250 by rigorous intervals. At $\Delta=-4075$ a direct complex determinant removes the sole remaining sextic. In degree six there are three smaller conjugates with moduli in $(2365,2366)$, $(2846,2847)$ and $(3701,3703)$, together with a principal conjugate larger than $10^{80}$ in modulus. If either of the first two exponent gaps is at least 12, diagonal dominance of the determinant expansion suffices. For first gaps 1 through 11, dividing by the principal root's highest power reduces to 121 nonzero complex minors, with a uniformly smaller error for every last gap. The small sieve, large sieve, boundary numbers, and all 121 minors were rerun. The archive includes the original checkers and expected outputs.\cite{prior}

